\subsubsection{Parametric Quantum Circuits (Ansätze)}\label{sec:parametric-quantum-circuits-ansatz}

\definition{Ansatz} (or \definition{Ansätze}) is a physical/mathematical term that refers to an \textbf{educated guess} for the \textbf{form of the solution}. Instead of searching among \textbf{all possible quantum states} (or in general, all possible solutions), we can restrict our search to a \textbf{subset of states} that are generated by a \textbf{parametric quantum circuit}, also called \textbf{ansatz}.

\highspace
In other words, an \important{Ansatz} is a \textbf{parametric quantum circuit}, meaning a circuit whose behavior depends on one or more adjustable parameters. Specifically, some quantum gates contain parameters that can be changed during optimization. However, since the quantum circuit is executed on a quantum computer, the \textbf{optimizer} (which is executed on a classical computer) \textbf{cannot modify the circuit structure}, but only the values of the parameters.

\highspace
\begin{flushleft}
    \textcolor{Green3}{\faIcon{question-circle} \textbf{What is parametrized in a parametric quantum circuit?}}
\end{flushleft}
Some gates, such as the Hadamard gate, are fixed and do not contain any parameters. On the other hand, gates like the rotation gates (e.g., $R_x(\theta)$, $R_y(\theta)$, $R_z(\theta)$, \autopageref{eq:rotation-operator-general}) have parameters (in this case, the angle $\theta$) that can be adjusted during the optimization process. The optimizer will modify these parameters to find the optimal solution to the problem at hand. This is extremely important because it means that with \textbf{one single quantum circuit} structure, we can \textbf{explore a wide range of quantum states by simply changing the parameters}, allowing us to efficiently search for the optimal solution \hl{without having to design a new circuit for each possible state.}

\highspace
\begin{flushleft}
    \textcolor{Green3}{\faIcon{question-circle} \textbf{What exits the ansatz?}}
\end{flushleft}
The output of an ansatz is not the solution itself, but rather a \textbf{quantum state candidate} that is generated by the parametric quantum circuit. 

\highspace
Let's suppose we have a Hamiltonian $H$ whose ground state is $\ket{\psi_{0}}$. The ansatz generates a quantum state $\ket{\psi(\theta)}$ that depends on the parameters $\theta$. The optimizer will minimize the expectation value $\bra{\psi(\theta)} H \ket{\psi(\theta)}$ with respect to the parameters $\theta$, until hopefully:
\begin{equation*}
    \ket{\psi(\theta^{*})} \approx \ket{\psi_{0}}
\end{equation*}
\textcolor{Red2}{\faIcon{exclamation-triangle}} However, there is \textbf{no guarantee} that the ansatz will be able to generate a state that is close to the ground state $\ket{\psi_{0}}$. An ansatz does \textbf{not solve the problem}. It \hl{defines the set of solutions that the optimizer is allowed to search through.}

\highspace
\begin{flushleft}
    \textcolor{Red2}{\faIcon{balance-scale} \textbf{Expressibility}}
\end{flushleft}
The \textbf{expressibility} of an ansatz is a measure of \hl{how well the ansatz can represent a wide range of quantum states}. Different ansätze have different expressibility, and different complexity. Intuitively, an ansatz with higher expressibility can represent a larger variety of quantum states, because it has a larger search space. However, a more expressive ansatz usually means a more complex quantum circuit, which can be more difficult to optimize, may require more quantum resources (e.g., more qubits, more gates, etc.), and may be more prone to noise and errors. Therefore, there is a \textbf{trade-off between expressibility and complexity} when designing an ansatz for a VQA.

\highspace
Since the ansatz will be \textbf{executed on a NISQ device}, it should be \textbf{shallow} (i.e., it should have a small depth) and must \textbf{use gates that are native to the hardware}. In general, if the ansatz expressibility is too low, it may not be able to represent the optimal solution, and the optimization process will fail to find a good solution. On the other hand, if the ansatz expressibility is too high, noise destroys the computation. Therefore, designing an ansatz with the right balance of expressibility and complexity is crucial for the success of a VQA.